\documentclass[11pt]{article}
\usepackage{amsmath,amssymb,hyperref}
\title{Gapped Reciprocal Covers and the Fixed Cube Obstruction}
\author{}
\date{November 22, 2026}
\begin{document}
\maketitle

\paragraph{Scope.}
This note proves a transfer of a geometric shattering construction to
positive affine reciprocals on a specially chosen finite point set.
It also proves that an affine or common-positive-denominator projective
embedding of the entire Boolean cube cannot provide the required
halfspace cover. It does not prove an upper bound on $W(n)$.

\paragraph{A gapped cover transfer.}
Let $P$ be a finite set, let $T\subseteq P$, and let $B_1,\ldots,B_k$
be strictly positive functions on $P$. Suppose every $p\in T$ has
some $B_h(p)\le s$, while $B_h(p)\ge L$ for every $p\notin T$ and
all $h$. If
\[
 0<s,\qquad 0<L,\qquad \frac{k}{L}<\tau<\frac1s,
\]
then
\[
 F(p)=\sum_{h=1}^k\frac1{B_h(p)}-\tau
\]
is strictly positive on $T$ and strictly negative on $P\setminus T$.
For a covered point one summand is at least $1/s$ and all others
are positive; for an uncovered point the sum is at most $k/L$.
This implication is formalized as
\texttt{HeadComplexity.reciprocal\_sum\_separates\_gapped\_cover}.

\paragraph{Application to geometric shattering.}
Theorem 1 and Proposition 5 of Csik\'os, Kupavskii, and Mustafa,
\emph{Tight Lower Bounds on the VC-dimension of Geometric Set Systems}
(JMLR 20, 2019), construct, for $d\ge4$ and $k\ge2$, a positive
point set of size $\Omega(dk\log k)$ such that each of its subsets is
a union of at most $k$ origin-anchored boxes. Their proof permits
coordinatewise order-preserving stretching so that consecutive
positive coordinate values have ratio greater than any prescribed
$R>1$. The combinatorial box construction comes from Lemma 2 and
Section 4 of Kupavskii, Mustafa, and Pach, \emph{New Lower Bounds for
$\epsilon$-Nets} (SoCG 2016).

Here is a direct inference for reciprocal scores. Stretch with
$R=4k(d+1)$. Each selected box can be written
$[0,b_{h1}]\times\cdots\times[0,b_{hd}]$, with each endpoint a
coordinate value of the stretched point set. Set
\[
 B_h(p)=1+\sum_{i=1}^d p_i/b_{hi},\qquad
 s=d+1,\qquad L=R,\qquad \tau=1/(2(d+1)).
\]
For a point inside a selected box, $B_h(p)\le d+1$.
For a point outside that box, some coordinate ratio exceeds $R$,
so $B_h(p)>R$. The preceding transfer applies. All denominator
constants and slopes are strictly positive and all numerators are
one. If fewer than $k$ boxes are used, the same inequalities hold
with their actual number. The empty target is represented by a
negative constant, without any heads.

Consequently there exist point sets of size $\Omega(dk\log k)$ in
$\mathbb R_{>0}^d$ shattered by at most $k$ positive affine reciprocal
terms and a constant readout. This concerns an existentially chosen
point set, not the prescribed $2^n$ Boolean vertices.

\paragraph{Parity prevents an affine transfer.}
Let $n\ge1$, and let $E$ be the even parity class of $\{0,1\}^n$.
Every strict halfspace whose intersection with the cube is contained
in $E$ contains at most one cube vertex. Indeed, suppose it contains
distinct $u,v\in E$, and write its affine defining function as $L$,
with $L(u),L(v)>0$. Choose a coordinate where $u$ and $v$ differ,
and swap that coordinate between them, producing $w,z$.
Both new vertices have odd parity, and $u+v=w+z$. Hence
\[
 0<L(u)+L(v)=L(w)+L(z)\le0,
\]
a contradiction. Toggling one fixed coordinate gives a bijection
between the two parity classes, so $|E|=2^{n-1}$.
Any union of strict halfspaces representing $E$ therefore has at
least $2^{n-1}$ members.

This lower bound persists after any affine feature map, since the
preimage of a halfspace is a halfspace. It also persists for a map
\[
 p_i(x)=A_i(x)/D(x),\qquad D(x)>0,
\]
with affine $A_i,D$ and a common denominator: multiplying a halfspace
inequality by $D$ gives an affine inequality in $x$. The scalar
identity is formalized as
\texttt{HeadComplexity.projective\_halfspace\_sign\_iff}.
This excludes those direct transfers of the geometric construction.
It is not a lower bound for sums of general affine fractions, which
can represent parity much more efficiently.

\paragraph{What is missing.}
The geometric construction uses non-affine stretching and a special
point set. No admissible one-layer compilation of a suitable
cube-compatible replacement is established. Merely asking for such
a replacement at the target count would be a theorem-strength gap,
not a proof mechanism.

\paragraph{Primary sources.}
\url{https://www.jmlr.org/papers/volume20/18-719/18-719.pdf}

\url{https://lipn.univ-paris13.fr/~mustafa/Research/Papers/epsnetslowerbounds.pdf}

\end{document}
